\begin{theorem}[节点坐标域上两条单生成整序的判别式与指数比值]\label{thm:xi-terminal-zm-delta-node-order-discriminant-index-ratio}
取 \(\overline{\mathcal C}_\delta\) 的任一有限节点 \((y_0,T_0)\)。
令
\[
K_0:=\QQ(y_0)=\QQ(T_0),
\qquad
\mathcal O_y:=\ZZ[y_0]\subset K_0,
\qquad
\mathcal O_T:=\ZZ[T_0]\subset K_0,
\]
并令 \(\mathcal O_{K_0}\) 为 \(K_0\) 的整数环。
则有严格恒等式
\[
\frac{[\mathcal O_{K_0}:\mathcal O_T]}{[\mathcal O_{K_0}:\mathcal O_y]}
=\frac{2^{7}\cdot 3\cdot 307}{31},
\qquad
\frac{\Disc(\mathcal O_T)}{\Disc(\mathcal O_y)}
=\left(\frac{2^{7}\cdot 3\cdot 307}{31}\right)^2.
\]
等价地，对任意素数 \(\ell\) 有逐素数赋值差分
\[
v_\ell([\mathcal O_{K_0}:\mathcal O_T]) - v_\ell([\mathcal O_{K_0}:\mathcal O_y])
=
v_\ell\!\left(\frac{2^{7}\cdot 3\cdot 307}{31}\right),
\]
从而
\[
v_2(\cdot)=7,\qquad v_3(\cdot)=1,\qquad v_{307}(\cdot)=1,\qquad v_{31}(\cdot)=-1,
\]
并且在 \(\ell\in\{11,727\}\) 处差分为 \(0\)。
\end{theorem}
\begin{proof}
由命题 \ref{prop:xi-terminal-zm-delta-nodes-elimination-a5}，节点坐标满足 \(\QQ(y_0)=\QQ(T_0)\)。
由于 \(Q_5(y)\) 与 \(A_5(T)\) 均为首一整系数多项式，\(y_0,T_0\) 均为代数整数，故 \(\mathcal O_y,\mathcal O_T\) 为 \(K_0\) 的整序。
\par
由结论 \ref{con:xi-terminal-zm-leyang-q5-discriminant-727} 与结论 \ref{con:xi-terminal-zm-delta-a5-discriminant-quadratic}，
\[
\Disc(\mathcal O_y)=\Disc(Q_5)
=-2^{28}\cdot 3^{9}\cdot 11^{2}\cdot 31^{3}\cdot 727,
\]
\[
\Disc(\mathcal O_T)=\Disc(A_5)
=-2^{42}\cdot 3^{11}\cdot 11^{2}\cdot 31\cdot 307^{2}\cdot 727.
\]
对同一数域 \(K_0\) 的任意整序 \(\mathcal O\subseteq \mathcal O_{K_0}\)，判别式满足
\(
\Disc(\mathcal O)=\Disc(\mathcal O_{K_0})\,[\mathcal O_{K_0}:\mathcal O]^2
\)。
将其分别应用于 \(\mathcal O_y,\mathcal O_T\) 并取比值得
\[
\frac{\Disc(\mathcal O_T)}{\Disc(\mathcal O_y)}
=\left(\frac{[\mathcal O_{K_0}:\mathcal O_T]}{[\mathcal O_{K_0}:\mathcal O_y]}\right)^2
=\frac{2^{42}\cdot 3^{11}\cdot 11^{2}\cdot 31\cdot 307^{2}\cdot 727}{2^{28}\cdot 3^{9}\cdot 11^{2}\cdot 31^{3}\cdot 727}
=\left(\frac{2^{7}\cdot 3\cdot 307}{31}\right)^2.
\]
取正平方根得到指数比值恒等式。
逐素数差分由对两侧取 \(v_\ell\) 并除以 \(2\) 得出。证毕。
\end{proof}
\leanverified{paper\_xi\_terminal\_zm\_delta\_node\_order\_discriminant\_index\_ratio}

\endinput
